% method_hmm.tex
% Subsection 3.1 body. Flat prose, no subsubsections. Topic flow
% (language-audit pass, 2026-07-14; see project_language_audit_2026-07-14
% memory): dataset/return definition -> tuple definition -> state
% partition -> emissions -> transition estimation + jump/tail override
% -> hyperparameter calibration. Split the original single dense
% tuple-plus-state-plus-emission paragraph into three, moved the
% variance/sensitivity caveats out of the model-definition flow, and
% cut a prose walk-through of the generator that duplicated Algorithm 1.
% The bridge to the multi-asset composition lives at the top of
% method_sim.tex, not here (avoids restating the per-ticker-marginal
% summary twice across the subsection boundary).

The dataset was $424$ United States-listed equities and
exchange-traded funds spanning ten years ($2014$--$2024$). To validate
the hybrid hidden Markov model on a single representative asset,
we focused the main analysis on the SPDR S\&P 500 ETF (SPY), reserving the $2,766$ observations from January 3,
$2014$ to December 31, $2024$ for in-sample fitting and the subsequent
$249$ trading days of $2025$ for out-of-sample testing. The data
consisted of daily open, high, low, and close prices and
volume-weighted average price values. We computed the excess growth rate $G_{i,j}$ for ticker $i$ between
trading days $j-1$ and $j$ from the equity price series
$P_{i,j}$ as:
\begin{equation}
    G_{i,j} \;\equiv\; \left(\frac{1}{\Delta t}\right)\,
    \ln\!\left( \frac{P_{i,j}}{P_{i,j-1}} \right) \;-\; r_f,
    \label{eq:growth_rate}
\end{equation}
where $\Delta t = 1/252$ is the fraction of a trading year spanned by
one trading day, and $r_f$ is a constant continuously compounded
risk-free rate derived from Treasury STRIPS (Separate Trading of
Registered Interest and Principal of Securities) bond yields. The excess growth rate
has units of $\mathrm{year}^{-1}$ and is time-additive under
continuous compounding.

We defined the hybrid hidden Markov model with a jump-duration
mechanism (HMM-WJ) as the tuple
$\mathcal{M} = (\mathcal{S}, \mathcal{O}, \mathbf{T}, \mathbf{E}, \bar{\pi})$~\cite{rabinerJuang1986}
(Figure~\ref{fig:model_architecture}). The hidden state
$S_t \in \mathcal{S} = \{1,\dots,N\}$ represents a market regime, and
the observation at step $t$ is the excess growth rate $G_t$ of
Eq.~\eqref{eq:growth_rate}, drawn from the continuous observation
space $\mathcal{O} \subset \mathbb{R}$ (we suppress the ticker index
$i$ when only one asset is in play). The transition matrix
$\mathbf{T}$ governs regime-to-regime transitions, $\mathbf{E}$ is the
state-conditional emission model, and $\bar{\pi}$ is the stationary
distribution.

States were defined by quantile boundaries of a fitted Laplace
cumulative distribution: $Q_k = F_L^{-1}(k/N;\, \mu_L, b_L)$, where
$F_L^{-1}(\cdot;\, \mu_L, b_L)$ is the inverse cumulative distribution
function of the Laplace distribution with location $\mu_L$ and scale
$b_L$ (both estimated from the in-sample growth-rate history by
maximum likelihood), with the endpoints fixed at $Q_0 = -\infty$ and
$Q_N = +\infty$. An observation was assigned to state $k$ when
$Q_{k-1} < G_t \le Q_k$. We chose the Laplace distribution for its
sharper peak, which better matches the concentration of small price
movements in the empirical marginal density
(Fig.~\ref{fig:empirical_motivation}(a))~\cite{kotzKozubowskiPodgorski2001, tothJones2019},
and for its closed-form quantile function, which scales to the
$424$-ticker universe without iterative optimization~\cite{bilmes1998}.

The within-state emission followed a location-scale Student-$t$
distribution:
\begin{equation}
    G_t \mid S_t = k \;\sim\; \mu_k + \sigma_k \cdot t_{\nu},
    \qquad \nu = 5,
    \label{eq:emission}
\end{equation}
where $t_\nu$ denotes a standard Student-$t$ random variable with
$\nu$ degrees of freedom, $\mu_k$ is the sample mean of the in-sample
observations assigned to state $k$, and $\sigma_k$ is their sample
standard deviation; because a standard $t_\nu$ has variance
$\nu/(\nu-2)$, $\sigma_k$ is a scale rather than the emission's actual
standard deviation, which is $\sigma_k\sqrt{5/3}$ at $\nu = 5$. We
selected $\nu = 5$ from a sensitivity sweep over
$\nu \in \{3, 4, 5, 6, 7, 8, 10, 15, 30, \infty\}$
($\nu = \infty$ recovers the Normal baseline) as the value that
minimized the kurtosis gap without degrading distributional fidelity;
the full comparison against a Gaussian-emission HSMM baseline of
Bulla and Bulla~\cite{bullaBulla2006} is given in
Table~\ref{tab:emission_hsmm} of
Sec.~\ref{sec:supp-emission-hsmm}.

With the state space and emission family fixed, we estimated the
transition matrix $\mathbf{T}$ by counting observed state-transition
frequencies directly, avoiding the initialization sensitivity of
iterative expectation-maximization~\cite{bilmes1998}. Direct
counting, however, gives empirical transition matrices low
probabilities for exiting central states into tail states, so models
built this way lose the effect of volatility clustering too
quickly~\cite{bullaBulla2006}. We corrected
this by adding two jump hyperparameters, $\epsilon$ representing the
probability of a jump event and $\lambda$ representing the mean
jump-event duration, along with a tail-state selection rule. Tail
states were defined as $\mathcal{S}_{-} = \{1, \dots, N_{\rm tail}\}$
and $\mathcal{S}_{+} = \{N - N_{\rm tail} + 1, \dots, N\}$, where
$N_{\rm tail}$ was user-configurable. When a jump was triggered, the
jump-episode length $K$ was sampled from a Poisson distribution
$K \sim \mathrm{Poisson}(\lambda)$~\cite{glasserman2003}; since the
Poisson support includes zero, a draw of $K = 0$ forces no tail steps
and the chain takes an ordinary transition, so $\lambda$ sets the mean
length of a triggered episode rather than a guaranteed minimum dwell.
During the $K$ forced steps, the standard Markovian transitions were
overridden to target tail states: at each forced step the tail set was
drawn as $\mathcal{S}_{-}$ with a user-configurable probability
$p_{\rm neg}$ (default $0.52$) and as $\mathcal{S}_{+}$ otherwise, a
bias toward negative-tail states that reproduced gain/loss asymmetry
(Figure~\ref{fig:model_architecture}). The empirical
transition matrix exhibited a dominant near-diagonal band reflecting
short-range regime persistence, with natural residence times of
$1$--$2$ steps in every state including the tails
(Fig.~\ref{fig:model_internals} in Online
Appendix~\ref{sec:supp-transitions}).
The stationary distribution $\bar{\pi}$ satisfies
$\bar{\pi} = \bar{\pi}\mathbf{T}$ and was estimated by matrix powering
($\mathbf{T}^{50}$)~\cite{hamilton2018regime, grinsteadSnell1997}; the
initial state was drawn from $S_0 \sim \bar{\pi}$. Because $\bar{\pi}$
is the invariant distribution of $\mathbf{T}$ alone, it only
approximates the invariant distribution of the jump-augmented process,
whose state also includes jump status and remaining episode length;
this introduces a small, unquantified transient rather than exact
stationarity. Algorithm~\ref{alg:hmmwj} gives the full simulator: at
each step the chain takes a standard transition with probability
$1 - \epsilon$ or, otherwise, enters a Poisson jump episode that
forces $K$ consecutive draws from a tail set before normal
transitions resume, sampling an observation
$G_t \sim \mathbf{E}_{S_t}$ at every step.

% --- Architecture figure ----------------------------------------------------
\begin{figure}[tp]
\centering
\resizebox{\textwidth}{!}{%
\begin{tikzpicture}[
  state/.style    = {circle, draw=black, very thick, minimum size=1.1cm,
                     inner sep=0pt, font=\small},
  tailbot/.style  = {state, fill=red!25,  draw=red!70!black},
  tailtop/.style  = {state, fill=blue!25, draw=blue!70!black},
  midstate/.style = {state, fill=white,   draw=black!70},
  markov/.style   = {<->, >=Stealth, thin, gray!80},
  trigger/.style  = {->, >=Stealth, thick, dashed, black!80},
  jumparc/.style  = {->, >=Stealth, very thick},
]
  \node[tailbot]              (s1)   at ( 0.0, 0) {$1$};
  \node[tailbot]              (s2)   at ( 2.0, 0) {$2$};
  \node[font=\Large]          (ld)   at ( 4.0, 0) {$\cdots$};
  \node[midstate]             (sk)   at ( 6.0, 0) {$k$};
  \node[font=\Large]          (rd)   at ( 8.0, 0) {$\cdots$};
  \node[tailtop]              (sNm1) at (10.0, 0) {$N\!-\!1$};
  \node[tailtop]              (sN)   at (12.0, 0) {$N$};
  \draw[markov] (s1)  -- (s2);
  \draw[markov] (s2)  -- (ld);
  \draw[markov] (ld)  -- (sk);
  \draw[markov] (sk)  -- (rd);
  \draw[markov] (rd)  -- (sNm1);
  \draw[markov] (sNm1)-- (sN);
  \draw[red!75!black, thick]
    plot[domain=-0.42:0.42, samples=40, smooth, variable=\t]
    ({ 0.0+\t}, {1.45 + 0.72*exp(-\t*\t/0.024)});
  \draw[red!25, thin] (-0.42,1.45) -- (0.42,1.45);
  \draw[red!75!black, thick]
    plot[domain=-0.42:0.42, samples=40, smooth, variable=\t]
    ({ 2.0+\t}, {1.45 + 0.72*exp(-\t*\t/0.024)});
  \draw[red!25, thin] (1.58,1.45) -- (2.42,1.45);
  \draw[black!55, thick]
    plot[domain=-0.42:0.42, samples=40, smooth, variable=\t]
    ({ 6.0+\t}, {1.45 + 0.72*exp(-\t*\t/0.024)});
  \draw[gray!40, thin] (5.58,1.45) -- (6.42,1.45);
  \draw[blue!75!black, thick]
    plot[domain=-0.42:0.42, samples=40, smooth, variable=\t]
    ({10.0+\t}, {1.45 + 0.72*exp(-\t*\t/0.024)});
  \draw[blue!25, thin] (9.58,1.45) -- (10.42,1.45);
  \draw[blue!75!black, thick]
    plot[domain=-0.42:0.42, samples=40, smooth, variable=\t]
    ({12.0+\t}, {1.45 + 0.72*exp(-\t*\t/0.024)});
  \draw[blue!25, thin] (11.58,1.45) -- (12.42,1.45);
  \node[font=\footnotesize\itshape, text=black!55]
    at (6.0, 2.7) {$\mu_k\!+\!\sigma_k\cdot t_5$ emission};
  \node[draw, rounded corners=5pt, fill=yellow!20, very thick,
        minimum width=2.8cm, minimum height=1.1cm,
        align=center, font=\small] (clock) at (6.0, 5.0)
        {Poisson clock\\[1pt]prob.~$\epsilon$};
  \draw[trigger] (sk.north) -- (clock.south)
    node[midway, right=3pt, font=\footnotesize] {jump triggered};
  \draw[jumparc, red!70!black]
    (clock.west)
    to[out=180, in=100]
    node[above, sloped, font=\footnotesize, red!80!black,
         pos=0.45, yshift=2pt]
      {$K\!\sim\!\mathrm{Pois}(\lambda)$ forced steps}
    (s1.north);
  \draw[jumparc, blue!70!black]
    (clock.east)
    to[out=0, in=80]
    node[above, sloped, font=\footnotesize, blue!80!black,
         pos=0.45, yshift=2pt]
      {$K\!\sim\!\mathrm{Pois}(\lambda)$ forced steps}
    (sN.north);
  \draw[decorate,
        decoration={brace, amplitude=6pt, mirror},
        red!70!black, thick]
    (-0.6, -0.65) -- (2.6, -0.65)
    node[midway, below=7pt, red!80!black, font=\small]
      {$\mathcal{S}_{-}$\ (bottom tail)};
  \draw[decorate,
        decoration={brace, amplitude=6pt, mirror},
        blue!70!black, thick]
    (9.4, -0.65) -- (12.6, -0.65)
    node[midway, below=7pt, blue!80!black, font=\small]
      {$\mathcal{S}_{+}$\ (top tail)};
  \draw[->, >=Stealth, gray!70, thin]
    (-1.0, -0.5) -- (13.2, -0.5)
    node[right, font=\small, text=black!70] {state index};
  \node[font=\footnotesize, text=gray!80, above right=2pt and 4pt of sk]
    {prob.~$1\!-\!\epsilon$};
\end{tikzpicture}%
}
\caption{\textbf{Architecture of the hybrid hidden Markov model with a
    jump-duration mechanism (HMM-WJ).} At each step the chain transitions
    according to the empirical transition matrix $\mathbf{T}$
    (probability $1 - \epsilon$, thin bidirectional arrows) or enters a
    Poisson jump episode (probability $\epsilon$, dashed upward arrow
    to the Poisson clock). During a jump episode each of the
    $K \sim \mathrm{Poisson}(\lambda)$ forced steps independently draws
    the bottom tail set $\mathcal{S}_{-}$ (red, lowest-valued states)
    with probability $p_{\rm neg}$ or the top tail set $\mathcal{S}_{+}$
    (blue, highest-valued states) otherwise, before normal Markovian
    evolution resumes. Small bell curves above
    each state depict the state-conditional location-scale Student-$t$
    ($\nu=5$) emission distribution. Tail sets are defined as the
    $N_{\rm tail}$ lowest- and highest-quantile states under the
    Laplace quantile partition.}
\label{fig:model_architecture}
\end{figure}

\begin{algorithm}[t]
\caption{HMM-WJ single asset marginal generator.}
\label{alg:hmmwj}
\begin{algorithmic}[1]
  \REQUIRE Fitted parameters
    $(\mathbf{T}, \mathbf{E}, \bar{\pi}, \epsilon, \lambda,
    p_{\rm neg}, N_{\rm tail})$; horizon $T$;
    tail sets $\mathcal{S}_{-} = \{1, \dots, N_{\rm tail}\}$ and
    $\mathcal{S}_{+} = \{N - N_{\rm tail} + 1, \dots, N\}$.
  \ENSURE State sequence $\{S_t\}_{t=0}^{T}$ and growth-rate path
    $\{G_t\}_{t=1}^{T}$.
  \STATE Draw $S_0 \sim \bar{\pi}$;\quad $k \gets 0$
    \hfill \COMMENT{$k$: remaining forced-step counter}
  \FOR{$t = 1, \dots, T$}
    \IF{$k > 0$}
      \STATE $\mathcal{J} \gets \mathcal{S}_{-}$ w.p.\ $p_{\rm neg}$, else $\mathcal{S}_{+}$;\quad
        $S_t \sim \mathrm{Uniform}(\mathcal{J})$;\quad $k \gets k - 1$
        \hfill \COMMENT{forced tail step; sign drawn per step}
    \ELSIF{$\mathrm{Uniform}(0,1) > \epsilon$}
      \STATE $S_t \sim \mathbf{T}_{S_{t-1}, \cdot}$
        \hfill \COMMENT{standard Markov transition}
    \ELSE
      \STATE $K \sim \mathrm{Poisson}(\lambda)$
        \hfill \COMMENT{trigger jump; $K$ is the episode length}
      \IF{$K = 0$}
        \STATE $S_t \sim \mathbf{T}_{S_{t-1}, \cdot}$
          \hfill \COMMENT{$K = 0$: no forced steps, standard transition}
      \ELSE
        \STATE $\mathcal{J} \gets \mathcal{S}_{-}$ w.p.\ $p_{\rm neg}$, else $\mathcal{S}_{+}$;\quad
          $S_t \sim \mathrm{Uniform}(\mathcal{J})$;\quad
          $k \gets \min(K - 1, T - t)$
          \hfill \COMMENT{first forced step; $K$ steps total, sign per step}
      \ENDIF
    \ENDIF
    \STATE Draw $G_t \sim \mu_{S_t} + \sigma_{S_t}\, t_{5}$
      \hfill \COMMENT{Eq.~\eqref{eq:emission}}
  \ENDFOR
\end{algorithmic}
\end{algorithm}

We calibrated the two jump hyperparameters,
the jump probability $\epsilon$ and mean episode length $\lambda$,
drawing on the jump-diffusion tradition of~\cite{merton1976, kou2002}
while adopting an HSMM-style state-duration parameterization~\cite{bullaBulla2006} in place of the instantaneous-jump formulation
of those classical models. As a discrete-time state model, HMM-WJ
lengthens dwell time in tail states through the jump-duration
override and includes no continuous-time diffusion component.
Calibration used a multi-objective grid search that minimized the
discrepancy between historical and simulated market signatures,
targeting two stylized facts~\cite{cont2001, rydenTerasvirtaAsbrink1998} visible in the
empirical SPY history of Fig.~\ref{fig:empirical_motivation}:
volatility clustering (Fig.~\ref{fig:empirical_motivation}(d)),
measured as the squared error between the observed and simulated
autocorrelation function (ACF) of absolute excess growth rates~\cite{cont2001} up to a lag of $L = 25$ trading days, the range where
most of the observed ACF variation occurs (post-fit evaluation in Fig.~\ref{fig:model_comparison} extends to lag $252$ for the
visual comparison), and leptokurtosis
(Fig.~\ref{fig:empirical_motivation}(a, b)), captured by a penalty
term on the global kurtosis~\cite{mandelbrot1963}:
\begin{equation}
    J(\epsilon, \lambda) =
    \sum_{\tau=1}^{L} \left( \mathrm{ACF}_{\rm obs}(\tau)
    - \overline{\mathrm{ACF}}_{\rm sim}(\tau) \right)^2
    + w_\kappa \left( \kappa_{\rm obs} - \overline{\kappa}_{\rm sim} \right)^2,
    \label{eq:grid_search}
\end{equation}
where $\mathrm{ACF}_{\rm obs}(\tau)$ and $\kappa_{\rm obs}$ are the
empirical absolute-growth autocorrelation at lag $\tau$ and the
global kurtosis, respectively (kurtosis denoted $\kappa$ to avoid
collision with the jump-duration variable $K$), and the simulated
counterparts $\overline{\mathrm{ACF}}_{\rm sim}(\tau)$ and
$\overline{\kappa}_{\rm sim}$ were averaged over the jump-active paths
among $200$ independent synthetic paths of $2{,}766$ trading days
simulated per grid point: a path that fires no jump reduces to the
no-jump generator and provides no information about
$(\epsilon, \lambda)$, so the objective is conditioned on paths where
at least one jump occurred~\cite{glasserman2003}. We set the kurtosis-penalty weight to
$w_\kappa = 0.20$ to balance tail behavior against the temporal ACF
term. The grid explored $\epsilon$ from
$10^{-4}$ to $2.5 \times 10^{-2}$ and $\lambda$ from $10$ to $160$ in
brute-force evaluation. The resulting model produced a per-ticker
marginal that reproduces heavy tails, regime persistence, and
volatility clustering.
